\documentclass[11pt]{article}
\usepackage[T1]{fontenc}
\usepackage{lmodern,amsmath,amssymb,amsthm,mathtools,longtable,array}
\usepackage[margin=1.05in]{geometry}
\usepackage[hidelinks]{hyperref}
\newtheorem{theorem}{Theorem}[section]
\newtheorem{lemma}[theorem]{Lemma}
\newtheorem{proposition}[theorem]{Proposition}
\newtheorem{corollary}[theorem]{Corollary}
\theoremstyle{remark}\newtheorem{remark}[theorem]{Remark}
\newcommand{\R}{\mathbb R}\newcommand{\C}{\mathbb C}
\newcommand{\norm}[1]{\left\lVert#1\right\rVert}
\newcommand{\code}[1]{{\small\nolinkurl{#1}}}
\DeclareMathOperator{\Rea}{Re}\DeclareMathOperator{\Ima}{Im}
\DeclareMathOperator{\Log}{Log}
\title{An attained eventual reality threshold\protect\\
for a growing quartic arithmetic transform}
\author{Project Stardust}
\date{Reviewed editorial draft, October 29, 2026}
\begin{document}
\maketitle
\begin{abstract}
For the growing quartic cosine transform of the arithmetic
theta kernel, we prove that the set of eventual split-real
polynomial starting times is a closed upper ray with finite
endpoint strictly greater than eight. The proof includes
uniform actual saddle estimates in both the moving transition
region and its far complement, a finite exact Laguerre
certificate, and a constructive polynomial approximation
argument. A single radius confines all nonreal zeros at times
at least one. Nonreal zeros approaching the endpoint have a
subsequence colliding at a positive real zero of finite
multiplicity at least two. All zeros are eventually real and
simple. The meaning of eventual start is kept separate from
pointwise-in-time reality, and no quadratic threshold or
Riemann-hypothesis implication is asserted. The final statements
are mapped to a frozen Lean development. A separate general divisor-product
argument retains complex zeros and supplies a common Gaussian integration
budget under an explicit zero-strip hypothesis; it gives no quartic
propagation theorem. Their novelty remains
unresolved.
\end{abstract}

\section{The actual transform and its eventual starting times}
For $u\geq0$ define the arithmetic kernel
\begin{equation}\label{eq:K}
 K(u)=\sum_{n=1}^{\infty}
 (2\pi^2n^4e^{9u}-3\pi n^2e^{5u})e^{-\pi n^2e^{4u}},
\end{equation}
and, for real $t$ and complex $z$, put
\begin{equation}\label{eq:H}
 H_t(z)=\int_0^\infty e^{tu^4}K(u)\cos(zu)\,du.
\end{equation}
The sign of the quartic tilt is positive. It is neither the
decaying quartic heat kernel nor the classical quadratic
deformation. Let $\mathsf{LP}(t)$ mean that real polynomials
splitting into real linear factors converge to $H_t$ locally
uniformly on the entire complex plane. Define
\begin{equation}\label{eq:starts}
 \mathcal S=\{T\in\R:\mathsf{LP}(t)\text{ for every }t\geq T\},
 \qquad A=\inf\mathcal S.
\end{equation}
The infimum becomes an ordinary finite real number once
nonemptiness and a lower bound are proved.

\begin{theorem}\label{thm:main}
The set $\mathcal S$ is nonempty and equals $[A,\infty)$ for a
finite $A>8$. For every $t\geq A$, all complex zeros of $H_t$
are real. For every $\varepsilon>0$ there is a time
$s\in(A-\varepsilon,A)$ at which $H_s$ has a nonreal zero.
There is one finite radius containing every nonreal zero at
every time $t\geq1$. Some sequence of nonreal zeros at times
below $A$ converges to a positive real zero of $H_A$ of finite
multiplicity at least two. For every sufficiently large time,
all zeros are real and simple.
\end{theorem}

The last statement is eventual simplicity, not simplicity at
every $t>A$. The eventual-start set in~\eqref{eq:starts} must
also be distinguished from $\{t:\mathsf{LP}(t)\}$. Equality
of the former with $[A,\infty)$ does not exclude earlier
individual reality times. No exact double multiplicity or
continuously tracked zero path is asserted.

The classical quadratic theory uses a real-zero propagation
property of a quadratic multiplier. De Bruijn's and Newman's
primary work supplies the relevant comparison, but it does not
permit that property to be silently applied to $e^{tu^4}$.
The whole-plane statement here requires actual quartic
estimates, including the region beyond every fixed frequency
window. It has no proved implication for the Riemann hypothesis.

\section{Arithmetic decay, entireness and time continuity}
All operations in~\eqref{eq:H} concern an actual Lebesgue
integral. We first verify its domains.

\begin{lemma}\label{lem:kernel}
The kernel is continuous and strictly positive on $[0,\infty)$.
There is $C>0$ such that
\begin{equation}\label{eq:decay}
 K(u)\leq C\exp\bigl(9u-\tfrac\pi2e^{4u}\bigr),\qquad u\geq0.
\end{equation}
For every integer $k\geq0$ and real $T,B$,
$u^ke^{Tu^4+Bu}K(u)$ is integrable on that half-line.
\end{lemma}
\begin{proof}
Each summand in~\eqref{eq:K} factors as
\[
 \pi n^2e^{5u}(2\pi n^2e^{4u}-3)e^{-\pi n^2e^{4u}}>0.
\]
For $v\geq1$ and $n\geq1$, $n^2v\geq(n^2+v)/2$. Hence
\[
 n^me^{-\pi n^2v}\leq n^me^{-\pi n^2/2}e^{-\pi v/2}.
\]
The first factors have a finite sum for $m=2,4$. Uniform
convergence on compact $u$-intervals proves continuity, and
the two bounds, with $e^{5u}\leq e^{9u}$, prove~\eqref{eq:decay}.
The positive first summand proves strict positivity.

Eventually $u^k\leq e^u$ and
$Tu^4+(B+11)u\leq(\pi/2)e^{4u}$. Substitution into
\eqref{eq:decay} bounds the tilted moment by $Ce^{-u}$.
Continuity handles its initial compact interval.
\end{proof}

\begin{proposition}\label{prop:entire}
The function $H_t$ is entire, even and conjugation compatible,
and $H_t(0)>0$. Its spatial derivatives are jointly continuous
in $(t,z)$. Time convergence is uniform on every compact
complex set. For each fixed real $t$ there is $C_t>0$ with
\begin{equation}\label{eq:growth}
 |H_t(z)|\leq C_t e^{|z|^{3/2}}.
\end{equation}
\end{proposition}
\begin{proof}
On a compact time--space set, the $k$th derivative of the
exponential integrand is dominated by
$u^ke^{Tu^4+Bu}K(u)$ for fixed finite $T,B$. Lemma
\ref{lem:kernel} permits differentiation and continuity under
the integral. The two exponentials defining cosine give the
entire derivatives. Its elementary symmetries give evenness
and conjugation, while positivity gives $H_t(0)>0$.
Joint continuity on a compact time interval times a compact
complex set is uniform continuity, proving the compact-time claim.

For $R,u\geq0$, $Ru\leq u^3+R\sqrt R$: split at
$u=\sqrt R$. Also $u^3\leq u^4+1$. Using
$|\cos(zu)|\leq e^{|z|u}$ with $R=|z|$ gives
\[
 |H_t(z)|\leq e^{1+R^{3/2}}
                    \int_0^\infty e^{(t+1)u^4}K(u)\,du.
\]
The last integral is finite and positive. This proves the
specific subquadratic bound~\eqref{eq:growth}.
\end{proof}

\section{Real zeros and actual split-polynomial limits}
For this even growth class, all-real zeros give actual
polynomial approximation with its normalization restored.
We give a finite-disk proof rather than assume a converse
to a theorem about limits of zeros.

\begin{lemma}\label{lem:polynomial}
Let $f$ be entire and even, with $f(0)\in\R\setminus\{0\}$,
and suppose $|f(z)|\leq Ce^{|z|^{3/2}}$. If every zero of $f$
is real, then $f$ is a locally uniform limit on $\C$ of real
polynomials splitting over $\R$.
\end{lemma}
\begin{proof}
Choose radii $R\to\infty$ with no zero on their boundary;
isolated zeros exclude only countably many radii. List the
finitely many zeros $a$ in $|z|<R$, with multiplicity, and set
\[
 B_R(z)=\prod_{|a|<R}\frac{R(z-a)}{R^2-\overline a z}.
\]
No $a$ is zero. The denominators are nonzero on the closed
disk, and each factor has modulus one on its boundary.
The quotient $h_R=f/B_R$ extends across all zeros to a
zero-free holomorphic function on the disk, continuous and
holomorphic near its boundary. The maximum modulus principle gives
\[
 |h_R(z)|\leq M_f(R)\leq Ce^{R^{3/2}},\qquad
 |h_R(0)|\geq|f(0)|,
\]
the last inequality using $|B_R(0)|=\prod|a|/R\leq1$.

The zero multiset occurs in pairs $a,-a$ with equal
multiplicities. Each paired factor is even, so $B_R$ and $h_R$
are even. Choose the analytic logarithm $g_R$ on the disk
satisfying $e^{g_R}=h_R/h_R(0)$ and $g_R(0)=0$.
Uniqueness of that normalized logarithm makes it even.
For some positive $A_R=O(1+R^{3/2})$ we have
$\Rea g_R\leq A_R$.

Here is the needed logarithm estimate explicitly. The function
\[
 v_R=\frac{g_R}{2A_R-g_R}
\]
is holomorphic and bounded in modulus by one, because
\[
 |2A_R-g_R|^2-|g_R|^2
                         =4A_R(A_R-\Rea g_R)\geq0.
\]
It is even and vanishes at zero. The Schwarz lemma applied
with this double zero gives $|v_R(z)|\leq(|z|/R)^2$.
Solving for $g_R$ yields
\[
 |g_R(z)|\leq\frac{2A_R(|z|/R)^2}{1-(|z|/R)^2}.
\]
On each fixed compact disk this tends to zero as $R\to\infty$.
Consequently $h_R/h_R(0)\to1$ locally uniformly.

We also need a zero count to control the denominator product.
Choose a zero-free boundary radius $\rho\in(2R,3R)$. Harmonicity
of $\log|h_\rho|$ and its boundary values give
\[
 \sum_{|a|<\rho}\log\frac{\rho}{|a|}
 \leq\log M_f(\rho)-\log|f(0)|.
\]
Indeed this follows by taking the boundary mean of
$\log|h_\rho|=\log|f|$ and using
$f(0)=B_\rho(0)h_\rho(0)$. Every zero in $|a|<R$ contributes
at least $\log2$, so its counting function satisfies
$n_f(R)=O(1+R^{3/2})$.

Now normalize the finite product:
\[
 \frac{B_R(z)}{B_R(0)}=\frac{N_R(z)}{D_R(z)},\qquad
 N_R(z)=\prod_{|a|<R}(1-z/a),\quad
 D_R(z)=\prod_{|a|<R}(1-\overline a z/R^2).
\]
Pairing $a,-a$ turns $D_R$ into a product of
$1-\overline a^{,2}z^2/R^4$. On $|z|\leq r$, the sum of the
absolute pair corrections is at most
$n_f(R)r^2/(2R^2)=O_r(R^{-1/2})$.
The finite-product inequality
$|\prod_j(1-w_j)-1|\leq\exp(\sum_j|w_j|)-1$ therefore gives
$D_R\to1$ locally uniformly. The exact identity
\[
 \frac f{f(0)}=\frac{B_R}{B_R(0)}\frac{h_R}{h_R(0)}
\]
implies $N_R\to f/f(0)$ locally uniformly. Since all $a$ are
real, $f(0)N_R$ are the required split-real polynomials.
\end{proof}

\begin{lemma}\label{lem:reverse}
A nonzero entire locally uniform limit of split-real
polynomials has only real zeros.
\end{lemma}
\begin{proof}
For every real root $r$, $|x+iy-r|$ is nondecreasing in $|y|$.
Multiplying its linear factors gives
$|P(x+ia)|\leq|P(x+ib)|$ whenever $a^2\leq b^2$.
The inequality survives the pointwise limits. A nonreal zero
$x+iy$ of the limit would force zeros at all
$x+iy/(n+1)$, accumulating at $x$. The identity theorem would
make the limit zero everywhere, a contradiction.
\end{proof}

\begin{corollary}\label{cor:LP}
For every real $t$, $\mathsf{LP}(t)$ holds if and only if every
zero of the actual $H_t$ is real.
\end{corollary}
\begin{proof}
Use Proposition~\ref{prop:entire} and the preceding two lemmas.
\end{proof}

\section{An actual arithmetic obstruction at time eight}
This section records a computer-assisted finite certificate,
with analytic errors kept separate from the finite checks.
Let $x_0=632083/10000$ and $a=8/5$. Let $K_4$ be the first
four terms of~\eqref{eq:K}, and put
\[
 c_0(u)=\cos(x_0u),\quad c_1(u)=-u\sin(x_0u),\quad
 c_2(u)=-u^2\cos(x_0u),\quad F_k(u)=e^{8u^4}K_4(u)c_k(u).
\]

\begin{lemma}\label{lem:cutoff}
For $k=0,1,2$,
\[
 \left|H_8^{(k)}(x_0)-\int_0^a F_k(u)\,du\right|\leq10^{-16}.
\]
\end{lemma}
\begin{proof}
On $[0,a]$ we have $u^k\leq3$. The positive exponential series
gives $e^{4u}\geq1+4u+(4u)^4/24$, hence
$3n^2(e^{4u}-1)\geq9u+8u^4$ for $n\geq1$.
Using $3<\pi<4$, the omitted positive kernel terms satisfy
\begin{align*}
 u^ke^{8u^4}(K-K_4)(u)
 &\leq96\sum_{n\geq5}n^4e^{-3n^2}
 \leq192\sum_{n\geq5}e^{-2n^2}\\
 &<384e^{-50}<384(2/5)^{50}<10^{-17}.
\end{align*}
Here $n^4\leq2e^{n^2}$ follows from the exponential series,
and $2(5+j)^2\geq50+20j$ gives the geometric sum bound.
Also $e>5/2$.

The constant in~\eqref{eq:decay} may be taken at most $176$:
for $m=2,4$,
\[
 \sum_{n\geq1}n^me^{-\pi n^2/2}
 \leq2\sum_{n\geq1}e^{-n^2/2}
 \leq\frac2{e^{1/2}-1}\leq4,
\]
and $2\pi^2\cdot4+3\pi\cdot4\leq176$.
For $u\geq a$, the functions $u^4e^{-4u}$ and $ue^{-4u}$
decrease. The rational inequality
\[
 8a^4+111a<\tfrac32(5/2)^6<\tfrac32e^{4a}
\]
therefore gives $8u^4+111u\leq(3/2)e^{4u}$ for all $u\geq a$.
Since $u^k\leq e^{2u}$, the full tail integrand is bounded by
$176e^{-100u}$. Its integral is at most
$(176/100)e^{-160}<(176/100)2^{-160}$.
Adding the two errors gives
$a10^{-17}+(176/100)2^{-160}<10^{-16}$.
\end{proof}

Divide $[0,a]$ into 160 cells, with centers $c_j=(2j+1)/200$
and half-width $h=1/200$. Let $T_{j,k}$ be the degree-15
Taylor polynomial of $F_k$ about $c_j$ and define
\[
 P_k=\sum_{j=0}^{159}\int_{c_j-h}^{c_j+h}T_{j,k}(u)\,du.
\]
Only its eight even coefficients contribute. The normalized
sixteenth derivative bound on each cell is
$|F_k^{(16)}(u)/16!|\leq10^{22}$. Consequently
\[
 \left|\int_0^aF_k-P_k\right|
 \leq160\frac{2\cdot10^{22}h^{17}}{17}
 \leq\frac{16}{10^{15}},\qquad
 |H_8^{(k)}(x_0)-P_k|\leq\frac{161}{10^{16}}.
\]
The frozen exact certificate checks the derivative bound on
all 160 whole cells and all 480 centered integral enclosures.
At common denominator $10^{22}$, its summed endpoint
numerators are
\[
\begin{array}{c|rr}
k&\text{lower}&\text{upper}\\\hline
0&-1641482419995&-1641482419749\\
1&3656317813&3656318043\\
2&-1460002952576&-1460002952347.
\end{array}
\]
The full cell table and an independent exact rational re-sum
are saved with the source dossier. The interval calculations
use outward dyadic rounding at scale $2^{192}$, 48-term
range-reduced exponential and trigonometric series, and the
normalized coefficient recursion
\[
 B_0=e^{\phi},\qquad
 B_n=\frac1n\sum_{j=1}^n jA_jB_{n-j},\quad
 A_j=\phi^{(j)}/j!,\quad B_j=(e^\phi)^{(j)}/j!.
\]
The phases $8u^4+bu-\pi n^2e^{4u}$ retain $b=9,5$ and
$n=1,\ldots,4$ separately. Every finite inclusion is proved
by kernel reduction in the audited Lean source; this is an
explicit computer-assisted step, not a floating-point probe.

\begin{proposition}\label{prop:obstruction}
At $x_0$ the actual transform satisfies
\begin{equation}\label{eq:negative}
 H_8'(x_0)^2-H_8(x_0)H_8''(x_0)<-2\cdot10^{-20}.
\end{equation}
There is an open neighborhood of $8$ in which every time has
a nonreal actual zero. Every eventual start is bounded below
by some fixed number strictly greater than $8$.
\end{proposition}
\begin{proof}
Enlarge the three table intervals by $161/10^{16}$. Their
channel-zero and channel-two upper bounds are negative.
Let their negatives be $b_0,b_2>0$ and let $b_1$ be the enlarged
absolute channel-one bound. Exact rational arithmetic gives
$b_1^2-b_0b_2<-2\cdot10^{-20}$, proving~\eqref{eq:negative}.

Every real split polynomial satisfies $P'^2-PP''\geq0$ on
the real axis. Away from its roots this follows by dividing
by $P^2$ and summing the positive terms
$m_r/(x-r)^2$; continuity covers its roots. Locally uniform
entire convergence also gives convergence of the first two
derivatives by the Cauchy formula. Thus~\eqref{eq:negative}
excludes $\mathsf{LP}(8)$. Joint derivative continuity
preserves the strict negative sign in an open time interval
about $8$, and Corollary~\ref{cor:LP} supplies its nonreal zeros.
A fixed bad time in that interval strictly above $8$ excludes
every eventual start below or equal to it.
\end{proof}

\section{The theta strip and its actual logarithmic derivative}
The large-frequency analysis takes place in
$\mathcal U=\{w:|\Ima w|<\pi/8\}$. The same series as in
\eqref{eq:K} defines a holomorphic kernel there. All constants
denoted $C,c>0$ in the next two sections are absolute unless
an explicit subscript records a fixed dependence. Their
values may change between inequalities. A uniform assertion
as $a\to\infty$ means that one cutoff in $a$ works for every
other parameter satisfying the displayed hypotheses.

\begin{lemma}\label{lem:strip}
The complex kernel is even and conjugation compatible on
$\mathcal U$. If $w=a+ib\in\mathcal U$, $a\geq0$, $L=e^{4a}$
and $V=\Rea e^{4w}$, then
\begin{equation}\label{eq:rho}
 K(w)=2\pi^2e^{9w-\pi e^{4w}}(1+\rho(w)),\qquad
 |\rho^{(j)}(w)|\leq C_j\bigl(e^{-4a}+(1+L)^je^{-3\pi V}\bigr)
\end{equation}
for $0\leq j\leq4$ when $V\geq1$.
For $s\geq0$ and $c_b=\cos(4b)>0$,
\begin{equation}\label{eq:compactkernel}
 |K(s+ib)|\leq Cc_b^{-3}e^{-s}.
\end{equation}
\end{lemma}
\begin{proof}
Normal convergence of the theta series and all its derivatives
holds on compact subsets where $\Rea e^{4w}>0$.
Conjugation follows termwise. To check evenness, write
$\vartheta(v)=\sum_{n\in\mathbb Z}e^{-\pi n^2v}$ and
$F(w)=e^w(\vartheta(e^{4w})-1)/2$. The Gaussian theta
identity $\vartheta(v)=v^{-1/2}\vartheta(1/v)$, with the
analytic right-half-plane square root, gives
$F(-w)=F(w)+\sinh w$ on $\mathcal U$.
This identity is the Gaussian Poisson summation formula;
equivalently it is the zero-argument modular theta identity
in~\cite{dlmf}. Direct differentiation gives
$(d^2/dw^2-1)F=8K$. Applying this operator to the reflection
identity proves $K(-w)=K(w)$.

Factor out the $n=1$ leading exponential. The remaining
factor is
\[
 1-\frac3{2\pi}e^{-4w}
 +\sum_{n\geq2}\left(n^4-\frac3{2\pi}n^2e^{-4w}\right)
                    e^{-\pi(n^2-1)e^{4w}}.
\]
Each of its first four derivatives is a finite linear
combination of terms bounded by a constant times
$e^{-4a}$ or
$(1+L)^j n^{4+2j}e^{-\pi(n^2-1)V}$.
For $V\geq1$, dividing the last sum by $e^{-3\pi V}$ leaves
a sum bounded by
$\sum_{n\geq2}n^{4+2j}e^{-\pi(n^2-4)}<\infty$.
This proves~\eqref{eq:rho}, including its derivative bounds.

For every $V>0$ and $m=2,4$,
$\sum_{n\geq1}n^me^{-\pi n^2V}\leq C_mV^{-(m+1)/2}$.
Indeed $n^me^{-\pi n^2V/2}\leq C_mV^{-m/2}$, and the
remaining decreasing Gaussian sum is bounded by its integral
from zero. Substitute $V=e^{4s}c_b$ in the two kernel terms.
They are bounded by
$Ce^{-s}(c_b^{-5/2}+c_b^{-3/2})$, which implies
\eqref{eq:compactkernel} since $0<c_b\leq1$.
\end{proof}

Write
\[
 \sigma_t(w)=4tw^3+\frac{K'(w)}{K(w)},\qquad
 W_{t,z}(w)=e^{tw^4+izw}K(w)
\]
where the denominator is nonzero. Under $V\geq c_0ta^3$,
$t\geq1$ and large $a$, formula~\eqref{eq:rho} gives
\begin{equation}\label{eq:slope}
 \sigma_t(w)=4tw^3+9-4\pi e^{4w}+\epsilon(w),\qquad
 |\epsilon^{(j)}(w)|\leq Ce^{-4a}\quad(0\leq j\leq3),
\end{equation}
for any fixed $c_0>0$. To see the uniformity, the other error
terms are at most $C\exp(16a-ca^3)$, which is eventually
bounded by $e^{-4a}$. The same bound holds along horizontal
shifts $w+s$, $s\geq-1$, after replacing the right side by
$Ce^{-4a}$: the potentially growing terms have the form
$\exp(16a+16s-c ta^3e^{4s})$, whose maximum for $s\geq-1$
occurs at the left endpoint once $a$ is large. In particular
the kernel is nonzero on those horizontal segments.

The distinction between the full slope $\sigma_t$ and its
leading model $4tw^3+9-4\pi e^{4w}$ remains essential.
Stationarity below always refers to the full slope.

\section{A uniform horizontal saddle estimate}
Fix $c_0=1/1000$. Call $(t,z,v)$ admissible when
\begin{equation}\label{eq:admissible}
 \begin{gathered}
 t\geq1,\quad v=a+ib\in\mathcal U,\quad a\text{ sufficiently large},
 \quad y=\Ima z\leq1,\\
 \sigma_t(v)+iz=0,\qquad
 \Rea e^{4v}\geq c_0ta^3.
 \end{gathered}
\end{equation}
The last condition is verified in the saddle constructions;
it prevents circular use of a theta asymptotic at an unknown
point. Set $Y=\max(0,-y)$, $Q=ta^3+Y$ and $L=e^{4a}$.

\begin{lemma}\label{lem:geometry}
Uniformly for admissible saddles,
\begin{equation}\label{eq:curvature}
 cQ\leq\Rea e^{4v}\leq CQ,\quad
 cL\leq|\beta|\leq CL,\quad
 cQ\leq\Rea\beta\leq CQ,\qquad
 \beta=-\sigma_t'(v)/2.
\end{equation}
For real $s\geq-1$,
\begin{equation}\label{eq:horizontalcurvature}
 \Rea\sigma_t'(v+s)\leq-cQ.
\end{equation}
On each sufficiently small fixed horizontal interval about
zero, $|\sigma_t''(v+s)|\leq CL$.
\end{lemma}
\begin{proof}
Take the real part of full stationarity in~\eqref{eq:slope}:
\begin{equation}\label{eq:stationaryreal}
 4\pi\Rea e^{4v}
       =4t(a^3-3ab^2)+9-y+\Rea\epsilon(v).
\end{equation}
Since $|b|<\pi/8<1/2$, $a$ is large and $y\leq1$, this is
comparable above and below to $ta^3+Y$. It also gives
$ta^3\leq CL$. In
\[
 \sigma_t'(v)=12tv^2-16\pi e^{4v}+\epsilon'(v),
\]
the first term has norm at most $CL/a$, whereas the
exponential term has norm $16\pi L$. This proves the norm
bounds for $\beta$. Taking real parts and using
\eqref{eq:stationaryreal} gives its real-part bounds, since
$ta^2=o(ta^3)$ uniformly.

For $s\geq-1$, $(a+s)^2\leq Ca^2e^{2s}$ and
$\Rea e^{4(v+s)}=e^{4s}\Rea e^{4v}$.
The positive polynomial part of
$\Rea\sigma_t'(v+s)$ is therefore at most
$(C/a)Qe^{2s}$. The negative exponential part is at most
$-cQe^{4s}$. Since $e^{2s}\leq e^2e^{4s}$ on this interval,
the negative term dominates for large $a$.
The uniform error~\eqref{eq:slope} is negligible as well,
proving~\eqref{eq:horizontalcurvature}.
Finally $\sigma_t''(w)=24tw-64\pi e^{4w}+\epsilon''(w)$;
the asserted small-window bound follows from $ta^3\leq CL$.
\end{proof}

For fixed admissible parameters, the shifted horizontal
integral is absolutely convergent: its strip margin is positive
and the double exponential dominates the quartic term at
infinity. Put
\[
 I_{t,z}(b)=\int_0^\infty W_{t,z}(s+ib)\,ds,\qquad
 G(\beta)=\int_\R e^{-\beta s^2}\,ds.
\]

\begin{proposition}\label{prop:horizontal}
There is a function $\eta(a)\to0$ such that, uniformly for
all admissible saddles,
\begin{equation}\label{eq:horizontal}
 \left|\frac{I_{t,z}(b)}{W_{t,z}(v)G(\beta)}-1\right|
                          \leq\eta(a).
\end{equation}
Also
\begin{equation}\label{eq:absolutemass}
 \int_0^\infty|W_{t,z}(s+ib)|\,ds
                     \leq C|W_{t,z}(v)|/\sqrt Q.
\end{equation}
\end{proposition}
\begin{proof}
On the horizontal interval $s\geq-1$, define
\[
 P(s)=\int_0^s(\sigma_t(v+r)+iz)\,dr.
\]
The elementary scalar ODE for the weight gives
$W_{t,z}(v+s)=W_{t,z}(v)e^{P(s)}$. This follows by
differentiating $W_{t,z}(v+s)e^{-P(s)}$, whose derivative
is zero; no complex neighborhood as wide as the horizontal
interval is assumed. Full stationarity gives $P'(0)=0$.
Lemma~\ref{lem:geometry} yields
\begin{equation}\label{eq:Pbounds}
 \Rea P(s)\leq-cQs^2\quad(s\geq-1),\qquad
 P(s)=-\beta s^2+O(L|s|^3)\qquad(|s|\leq\rho),
\end{equation}
where $\rho>0$ is a sufficiently small absolute constant.
The second estimate is real-parameter Taylor's theorem with
the uniform horizontal third-derivative bound.

For $\Rea\beta>0$, the complex Gaussian formula gives
$G(\beta)=\sqrt\pi/\sqrt\beta$, with its right-half-plane
branch. One direct verification differentiates under the
absolutely dominated integral and integrates the derivative
of $s e^{-\beta s^2}$: it obtains $G'=-G/(2\beta)$.
Thus $\sqrt\beta\,G$ is constant on the right half-plane,
and its value on the positive real axis is $\sqrt\pi$.
Hence $|G(\beta)|\asymp L^{-1/2}$ by
\eqref{eq:curvature}. Its two tails beyond $r>0$ satisfy
\begin{equation}\label{eq:oscgaussian}
 \left|\int_r^\infty e^{-\beta s^2}\,ds\right|
 +\left|\int_{-\infty}^{-r} e^{-\beta s^2}\,ds\right|
                           \leq\frac2{|\beta|r}.
\end{equation}
Integration by parts gives, for the positive tail,
$e^{-\beta r^2}/(2\beta r)
 -(2\beta)^{-1}\int_r^\infty e^{-\beta s^2}s^{-2}\,ds$;
taking absolute values proves the bound, and the negative tail
is identical. This estimate uses $|\beta|$, not a false
comparison of its real part with its modulus.

Take $r=L^{-7/16}$. On $[-r,r]$ the remainder in
\eqref{eq:Pbounds} is eventually at most one. The inequality
$|e^w-1|\leq e^{|w|}|w|$ bounds the integral difference
from $e^{-\beta s^2}$ by $CLr^4$. Relative to $G$, this is
$O(L^{3/2}r^4)=O(L^{-1/4})$. The Gaussian tails
\eqref{eq:oscgaussian}, after normalization, are
$O(L^{-1/2}r^{-1})=O(L^{-1/16})$.

It remains to estimate the actual annuli. For
$r\leq|s|\leq\rho$, reduce $\rho$ if necessary so that
\[
 |P'(s)|\geq cL|s|,\qquad |P''(s)|\leq CL.
\]
Indeed $P'(s)=-2\beta s+O(Ls^2)$ and
$|\beta|\geq cL$. On either annulus integrate
$(e^P)'/P'$ by parts. The boundary terms are bounded by
$C/(Lr)$. The integral of $|e^P P''/(P')^2|$ is at most
\[
 C\int_r^\rho\frac{ds}{Ls^2}\leq\frac{C}{Lr}.
\]
We used $|e^P|\leq1$ from~\eqref{eq:Pbounds}.
The relative annular error is again $O(L^{-1/16})$.

The remaining interval $[-1,-\rho]$ and right tail
$[\rho,\infty)$ are bounded absolutely by
$Ce^{-cQ\rho^2}$ and $Ce^{-cQ\rho^2}/(Q\rho)$ respectively.
After multiplication by $\sqrt L=e^{2a}$, both tend to zero
uniformly since $Q\geq a^3$.

Finally consider the physical prefix $0\leq s\leq a-1$.
Write $z=X+iy$. Equation~\eqref{eq:stationaryreal} gives
$c_b=\Rea e^{4v}/L\geq cQ/L$. Thus
\eqref{eq:compactkernel} and $-y s\leq Ys$ give
\[
 |W_{t,z}(s+ib)|
 \leq C e^{12a}Q^{-3}
          e^{t(a-1)^4+Ct+Y(a-1)-Xb}.
\]
The real quartic term on the line also contains
$-6ts^2b^2+tb^4\leq Ct$, which can be included in the exponent.
At the center,~\eqref{eq:rho} and
\eqref{eq:stationaryreal} give
\[
 |W_{t,z}(v)|\geq c
 e^{ta^4-ta^3-Cta^2+9a+Y(a-1/4)-Xb-Ca}.
\]
For $0<y\leq1$ the replacement of $-ya$ by this lower bound
costs only $Ca$. Subtract the two exponents. The polynomial
$t((a-1)^4-a^4+a^3+Ca^2+C)$ is at most $-ct a^3$
for large $a$, and the $Y$ term is $-3Y/4$. Therefore
\begin{equation}\label{eq:prefixbound}
 \frac{\int_0^{a-1}|W_{t,z}(s+ib)|\,ds}{|W_{t,z}(v)|}
                       \leq Ca\,e^{-cQ+Ca}.
\end{equation}
It too vanishes after Gaussian normalization.
The central, annular and three complementary bounds prove
\eqref{eq:horizontal}. One may take a constant multiple of
$L^{-1/4}+L^{-1/16}+(a+1)e^{-ca^3+Ca}$ for $\eta$,
after enlarging the constants. Absolute integration of
$e^{-cQs^2}$ on $[-1,\infty)$, together with
\eqref{eq:prefixbound}, proves~\eqref{eq:absolutemass}.
\end{proof}

This proof retains the actual horizontal slope and curvature.
The formal source often compares them to their leading model,
carrying its small linear mismatch explicitly. The present
grouping uses the same arithmetic estimates at the actual
saddle and proves the resulting bound directly. It does not
attribute this regrouping to an additional named Lean theorem.

\section{Signed contours, analytic pairing and exact local zeros}
The horizontal estimate concerns one half-line contribution.
We now relate it to the actual cosine integral, retaining
the endpoint terms when the two heights differ.

\begin{lemma}\label{lem:contour}
For real $h,k$ with $|h|,|k|<\pi/8$,
\begin{align}\label{eq:unequal}
 H_t(z)=\tfrac12\bigg(&I_{t,z}(h)
       +\overline{I_{t,\overline z}(k)}\notag\\
       &+i\int_k^hW_{t,z}(iy)\,dy\bigg).
\end{align}
In particular, choosing mirrored heights for the original
$z,-z$ contributions gives the exact identity
\begin{equation}\label{eq:mirror}
 H_t(z)=\tfrac12\bigl(I_{t,z}(h)+I_{t,-z}(-h)\bigr).
\end{equation}
\end{lemma}
\begin{proof}
Integrate $W_{t,z}$ around the rectangle with vertices
$0,R,R+ih,ih$. It is holomorphic on a neighborhood of this
rectangle in the theta strip. Its right vertical edge tends
to zero as $R\to\infty$: throughout that edge
$\Rea e^{4(R+iy)}\geq e^{4R}\min_{|v|\leq|h|}\cos(4v)>0$,
and the theta decay dominates every quartic and linear factor.
The two horizontal integrals are absolutely convergent for
the same reason. Their orientations therefore give
\[
 J_t(z):=\int_0^\infty W_{t,z}(u)\,du
       =I_{t,z}(h)+i\int_0^hW_{t,z}(iy)\,dy.
\]
Apply this also to $J_t(-z)$ at height $-k$.
Kernel evenness gives $W_{t,-z}(-iy)=W_{t,z}(iy)$, so its
vertical term is the negative integral from $0$ to $k$.
Conjugation gives
$I_{t,-z}(-k)=\overline{I_{t,\overline z}(k)}$.
Adding and dividing by two proves~\eqref{eq:unequal}.
For $k=h$ the vertical integral is zero, giving
\eqref{eq:mirror} in the original opposite-frequency notation.
\end{proof}

\begin{lemma}\label{lem:reflection}
For an admissible saddle with $y=\Ima z\leq0$,
\begin{equation}\label{eq:reflection}
 \frac{|I_{t,-z}(-b)|}{|W_{t,z}(v)G(\beta)|}
 \leq C\sqrt{L/Q}\,e^{2y(a-1)}+o(1),
\end{equation}
uniformly as $a\to\infty$. Consequently no zero can occur
at such an admissible frequency with $y\leq-1$, for large $a$.
\end{lemma}
\begin{proof}
At the same physical horizontal coordinate $s\geq0$,
conjugation of the kernel gives the exact modulus ratio
\[
 |W_{t,-z}(s-ib)|=|W_{t,z}(s+ib)|e^{2ys}.
\]
On $s\geq a-1$, use $e^{2ys}\leq e^{2y(a-1)}$ and
\eqref{eq:absolutemass}. On $0\leq s\leq a-1$, use
$e^{2ys}\leq1$ and~\eqref{eq:prefixbound}. Division by
$|G|\asymp L^{-1/2}$ proves~\eqref{eq:reflection}.
If $y\leq-1$, then
$\sqrt L\,e^{2y(a-1)}\leq e^{2a}e^{-2(a-1)}=e^2$.
Since $Q\geq a^3$, the right side tends to zero. The direct
contribution in~\eqref{eq:mirror}, normalized by the same
nonzero amplitude, tends to one by
Proposition~\ref{prop:horizontal}. Their sum cannot vanish.
\end{proof}

The remaining strip needs exact zero control. A real
conjugation-compatible function close to cosine can be shown
injective on the small disk containing a proposed conjugate
zero pair. The buffer in the next lemma is what permits the
derivative estimate.

\begin{lemma}\label{lem:smallcosine}
Let $a\in\R$ and let $f$ be holomorphic near the closed disk
$|q-a|\leq1$, with $f(\overline q)=\overline{f(q)}$ there.
If $|f(q)-\cos q|\leq\delta=1/1000$ on that disk, every zero
in $|q-a|\leq\delta$ is real and simple.
\end{lemma}
\begin{proof}
At a proposed zero $q_0$, $|\cos q_0|\leq\delta$.
Since $|\sin q|\leq e^{|\Ima q|}<3$ on the unit disk,
integration along the segment gives
$|\cos a|\leq\delta+3\delta=1/250$. Thus $|\sin a|\geq1/2$.
Set $d=-\sin a$. For $|q-a|\leq\delta$, the circle of
radius $1/2$ about $q$ stays inside the unit disk.
Cauchy's derivative estimate applied to $f-\cos$ gives
$|f'(q)+\sin q|\leq2\delta$.
Also $|\sin q-\sin a|\leq3\delta$, so
$|f'(q)-d|\leq5\delta\leq|d|/2$ on the small disk.

Integrate $f'-d$ on the segment between any two points of
that convex disk. The reverse triangle inequality gives
\[
 |f(q)-f(p)|\geq\tfrac12|d|\,|q-p|.
\]
The conjugate of $q_0$ is another zero in the disk, so the
inequality forces $q_0=\overline q_0$. The derivative bound
also gives $|f'(q_0)|\geq|d|/2>0$.
\end{proof}

\begin{lemma}\label{lem:straighten}
Let $g,E,\Phi$ be holomorphic on $|z-x|<r$, where $x\in\R$,
with $E$ nonzero and all three functions conjugation compatible.
Suppose $s>0$, $D\geq0$,
\[
 D/s\leq10^{-6},\quad 1/s\leq r/2,\quad
 |\Phi'(z)-s|\leq D,
\]
and
\begin{equation}\label{eq:weightedcosine}
 |g(z)/E(z)-\cos\Phi(z)|\leq10^{-6}\cosh(\Ima\Phi(z)).
\end{equation}
Every zero $z$ with $|z-x|<r/2$ is real and simple.
\end{lemma}
\begin{proof}
Put $p=\Rea z$, $a=\Phi(p)\in\R$,
$\ell(q)=p+(q-a)/s$ and $f(q)=g(\ell(q))/E(\ell(q))$.
The closed unit disk about $a$ maps strictly inside the
original disk, by the buffer and $|p-x|<r/2$.
Integration of $\Phi'-s$ gives
$|\Phi(\ell(q))-q|\leq D/s\leq\varepsilon=10^{-6}$.
Both arguments are in the disk of radius two about $a$;
there $|\sin|<9$ and $\cosh(\Ima\Phi)\leq9$.
Thus~\eqref{eq:weightedcosine} implies
$|f(q)-\cos q|\leq18\varepsilon<1/1000$.

At the zero, writing $\Phi(z)=u+iv$ gives
$|\cos(u+iv)|^2=\cos^2u+\sinh^2v$.
Equation~\eqref{eq:weightedcosine} therefore implies
$|\tanh v|\leq\varepsilon$, hence
$|v|\leq\operatorname{arctanh}\varepsilon\leq2\varepsilon$.
On the vertical segment from $p$ to $z$,
$\Rea\Phi'\geq s-D\geq s/2$. Since $\Phi(p)$ is real,
integration gives $|\Ima\Phi(z)|\geq s|\Ima z|/2$.
Consequently $q_0=a+s(z-p)$ is a zero of $f$ with
$|q_0-a|=s|\Ima z|\leq4\varepsilon<1/1000$.
Lemma~\ref{lem:smallcosine} makes $q_0$ real and simple.
Thus $z$ is real, and
$f'(q_0)=g'(z)/(sE(z))\ne0$ proves its simplicity.
\end{proof}

\begin{proposition}\label{prop:localactual}
Suppose a real frequency $x$ has an admissible actual saddle
$v_0=a+ib$, with large $a$. Then every zero $z=x+iy$ of
$H_t$ with $|y|<1$ is real and simple.
\end{proposition}
\begin{proof}
We first construct the local analytic quantities rather than
differentiate a varying principal argument. Lemma
\ref{lem:geometry} gives $|\sigma_t'(v_0)|\asymp L$.
The distance of $v_0$ from the theta-strip boundary is at
least $cQ/L$: indeed $\cos(4b)\geq cQ/L$ and the distance
is comparable to this cosine near either boundary.
Since $Q\geq a^3$, a disk of radius $C/L$ lies inside the
strip for any fixed $C$ and large $a$.
On this disk $|\sigma_t''|\leq CL$, by the differentiated
series estimate~\eqref{eq:rho}; the real-part margin changes
by only a bounded amount.

For $|z-x|<2$, solve
$\sigma_t(v)+iz=0$ by the contraction
\[
 v\longmapsto v-
       \frac{\sigma_t(v)+iz}{\sigma_t'(v_0)}
\]
on a disk $|v-v_0|\leq C/L$ with a sufficiently large fixed
$C$. Its center displacement is $O(1/L)$, and its derivative
is $O(1/L)$ there. It therefore maps the disk into itself and
is contractive for large $a$. The fixed point $v(z)$ depends
holomorphically on $z$, either by uniform iteration on smaller
frequency disks or the analytic inverse theorem, and
\begin{equation}\label{eq:branchderiv}
 v(z)=v_0+O(1/L),\qquad v'(z)=-i/\sigma_t'(v(z))=O(1/L).
\end{equation}
For $|z-x|\leq1$ its imaginary frequency is at most one.
The initial theta margin changes by only $O(1)$ and so first
remains at least $(c_0/2)ta^3$, enough to use the series
estimate~\eqref{eq:slope}. Taking real parts of its exact
stationarity then gives~\eqref{eq:stationaryreal} and restores
the larger lower bound $\Rea e^{4v(z)}\geq c_0t(\Rea v(z))^3$
for large $a$. Thus these saddles are admissible with the
original fixed constant.

The actual amplitude
\[
 \mathcal A(z)=W_{t,z}(v(z))\sqrt{\pi/\beta(z)},\qquad
 \beta(z)=-\sigma_t'(v(z))/2,
\]
is nonzero and holomorphic. The square root is fixed by
$\Rea\beta>0$. Stationarity cancels the derivative of the
center in the exponential weight, giving
\begin{equation}\label{eq:amplitudederiv}
 \frac{\mathcal A'}{\mathcal A}
     =iv(z)-\frac{\beta'(z)}{2\beta(z)}
     =iv(z)+O(1/L).
\end{equation}
Here $\beta'=O(1)$ follows from
$|\sigma_t''|\leq CL$ and~\eqref{eq:branchderiv}.

Choose an analytic $P$ with $P(x)=0$ and
$\mathcal A(z)=\mathcal A(x)e^{P(z)}$. Put
$P^*(z)=\overline{P(\overline z)}$, $a_0=\mathcal A(x)$ and
\[
 E(z)=|a_0|e^{(P(z)+P^*(z))/2},\qquad
 \Phi(z)=\arg a_0+\frac{P(z)-P^*(z)}{2i}.
\]
These are holomorphic and conjugation compatible; $E$ is
nonzero. The argument of $a_0$ is a fixed real number.
Direct expansion gives the exact identities
\begin{align}\label{eq:pairamplitudes}
 \tfrac12(\mathcal A(z)+\overline{\mathcal A(\overline z)})
          &=E(z)\cos\Phi(z),\notag\\
 |\mathcal A(z)|+|\mathcal A(\overline z)|
          &=2|E(z)|\cosh(\Ima\Phi(z)).
\end{align}
Equations~\eqref{eq:branchderiv}--\eqref{eq:amplitudederiv}
also give $|\Phi'(z)-a|\leq C/L$ on the unit frequency disk.

Apply~\eqref{eq:unequal} with heights
$h=\Ima v(z)$ and $k=\Ima v(\overline z)$.
The two horizontal errors are at most
$o(1)(|\mathcal A(z)|+|\mathcal A(\overline z)|)$, uniformly.
We must still bound the vertical residual. Its height gap is
$O(|y|/L)$ by~\eqref{eq:branchderiv}.
Throughout the interval between the heights, the cosine margin
is at least $cQ/L$: cosine is concave on $(-\pi/2,\pi/2)$,
so its minimum on that interval is at an endpoint.
Equation~\eqref{eq:compactkernel} at real coordinate zero
bounds the kernel there by $CL^3/Q^3$. Also
$|\Rea z|\leq CL$, as follows by taking norms in full
stationarity at the center. Thus the vertical residual is at most
\[
 \exp(Ca+t-Xh+C),\qquad X=\Rea z,
\]
after absorbing its polynomial factors and using
$|X||h-k|=O(1)$.
On the other hand~\eqref{eq:stationaryreal} at $v(z)$, with
$|y|\leq1$, gives
\[
 |\mathcal A(z)|\geq
 \exp\bigl(ta^4-Cta^3-Ca-Xh\bigr).
\]
Changing its real center by $O(1/L)$ changes only the stated
constants. The ratio of the residual to this amplitude is
at most $\exp(-ct a^4+Ca)$ for large $a$, hence tends to zero
uniformly. It cannot be discarded as an exact cancellation,
but it is negligible in this relative estimate.

By~\eqref{eq:pairamplitudes}, we have now proved
\[
 |H_t(z)/E(z)-\cos\Phi(z)|
       \leq o(1)\cosh(\Ima\Phi(z)),\qquad |z-x|\leq1.
\]
For large $a$ replace $o(1)$ by $10^{-6}$.
If a proposed zero has $|y|<1$, this inequality first gives
$|\Ima\Phi(z)|\leq2\cdot10^{-6}$ as in the preceding lemma.
Vertical integration of $\Phi'$ gives
$|\Ima\Phi(z)|\geq(a-C/L)|y|$, so actually $|y|<1/2$.
Apply Lemma~\ref{lem:straighten} on the unit disk, with
$s=a$, $D=C/L$ and $r=1$. Its buffer and scaled derivative
conditions hold for large $a$. The zero is real and simple.
\end{proof}

\section{Actual saddles in the transition and far regions}
The preceding estimates do not assume that a saddle exists.
We now supply that input in two overlapping frequency regimes.
The distinction is quantitative: a fixed multiple of the
moving real scale, and the region beyond that multiple.

\begin{lemma}\label{lem:realcenter}
For every sufficiently large $t$ there is a positive real
stationary center $u=u_t$ with $u_t\to\infty$ as $t\to\infty$.
Writing $B=4\pi e^{4u}$, it satisfies
\begin{equation}\label{eq:realscale}
 \sigma_t(u)=0,\qquad B=4tu^3+9+O(e^{-4u})\asymp tu^3.
\end{equation}
\end{lemma}
\begin{proof}
The positive continuous function $e^{tu^4}K(u)$ tends to zero
at infinity and has a global maximum on $[0,\infty)$.
For each fixed $R\geq0$, its maximum on $[0,R]$ is at most
$e^{tR^4}\max_{[0,R]}K$ for $t\geq0$. Its value at $R+1$
is $e^{t(R+1)^4}K(R+1)$, which eventually exceeds that bound.
Thus every maximizing center escapes to infinity.
At a positive center its ordinary derivative vanishes, so
$4tu^3+K'(u)/K(u)=0$. The real-axis instance of
\eqref{eq:rho} gives
$K'(u)/K(u)=9-4\pi e^{4u}+O(e^{-4u})$.
This uses only the real kernel asymptotic, independent of $t$,
and proves~\eqref{eq:realscale} without assuming it in advance.
\end{proof}

\begin{proposition}\label{prop:transition}
Fix $M\geq0$. For all sufficiently large $t$, every zero in
$|z|\leq MB$ is real and simple, with the same $u,B$ as in
Lemma~\ref{lem:realcenter}.
\end{proposition}
\begin{proof}
We construct saddles for all $|z|\leq MB$, $\Ima z\leq0$.
The anchor
\[
 q_z=u+\tfrac14\Log(1+iz/B)
\]
has $|q_z-u|\leq C_M$ and a fixed positive distance, depending
only on $M$, from the boundary of $\mathcal U$. Indeed
$\Rea(1+iz/B)\geq1$ and $|1+iz/B|\leq1+M$.
On the disk $|w-q_z|\leq C_M/u$, all kernel errors and their
derivatives in~\eqref{eq:slope} are $O_M(e^{-4u})$.
Equation~\eqref{eq:realscale} makes actual stationarity
equivalent to the fixed-point equation
\begin{equation}\label{eq:transitionfixed}
 w=u+\tfrac14\Log\left(1+\frac{iz}{B}
       +\frac{4t(w^3-u^3)+\epsilon(w)-\epsilon(u)}B\right).
\end{equation}
The expression inside the logarithm remains in a right
half-plane. Its perturbation from $1+iz/B$ is $O_M(1/u)$,
since $t/B=O(u^{-3})$. At the center the right side differs
from $q_z$ by $O_M(1/u)$, and its derivative in $w$ is
$O_M(tu^2/B)=O_M(1/u)$.
Choose the fixed constant in the disk radius large enough
for the center bound and then take $u$ large. The map is a
contraction of that disk into itself, giving an actual root
$v(z)$ with
\[
 \Rea v=u+O_M(1),\qquad
 \Rea e^{4v}\asymp_M B,\qquad
 e^{4\Rea v}\asymp_M B.
\]
In particular the theta margin is positive. Taking real parts
of stationarity, as in~\eqref{eq:stationaryreal}, improves
this to $\Rea e^{4v}\geq c_0t(\Rea v)^3$ for large $u$.
These saddles are admissible, and their ratio $L/Q$ is bounded
by a constant depending only on $M$.

Use Lemma~\ref{lem:reflection} on any lower-half-plane
frequency with $y\leq-A_M/u$. Since
$\Rea v-1\geq u/2$ for large $u$, its reflected relative
amplitude is at most $C_Me^{-A_M}+o(1)$.
Choose the fixed $A_M$ so this is below $1/4$; the direct
relative amplitude is within $1/4$ of one for large $u$.
Equation~\eqref{eq:mirror} excludes such zeros.
Conjugation excludes the corresponding upper region.
Thus every zero in the transition disk has
$|\Ima z|<A_M/u<1$.

At its real projection $x=\Rea z$ we have $|x|\leq MB$;
the same construction supplies a real-frequency admissible
saddle. Proposition~\ref{prop:localactual} now proves exact
reality and simplicity of the zero.
All cutoffs in this argument depend only on the fixed $M$,
not on the sampled frequency.
\end{proof}

\begin{lemma}\label{lem:farsaddle}
There is a sufficiently large absolute $C_0$ with the following
property. Let $t\geq1$, $\Ima z\leq0$, $X=|z|$ and
\[
 \alpha=\tfrac14\log(X/(4\pi)).
\]
If $\alpha$ is sufficiently large and
$X\geq C_0t(\alpha+1)^3$, there is an admissible actual saddle
$v$ with $\Rea v=\alpha+O(1)$, uniformly in these parameters.
\end{lemma}
\begin{proof}
First solve the leading model, which is entire and has no
theta-strip domain restriction. The anchor
$q=\tfrac14\Log(iz/(4\pi))$ has real part $\alpha$ and
imaginary part in $[-\pi/8,\pi/8]$, since $\Rea(iz)\geq0$.
Consider
\[
 w\longmapsto\tfrac14\Log\left(\frac{iz+4tw^3+9}{4\pi}\right)
\]
on a disk about $q$ of radius $Ct(\alpha+1)^3/X$.
Choose $C_0$ sufficiently large that this radius is below
$1/100$. On this disk $\Rea(w^3)\geq\alpha^3/2$ for large
$\alpha$, so the logarithm argument has positive real part.
Its modulus is comparable to $X$, since the cubic perturbation
has norm at most $CX/C_0$.
The center displacement is at most $Ct(\alpha+1)^3/X$ by
$|\Log(1+\zeta)|\leq2|\zeta|$ for $|\zeta|\leq1/2$;
its derivative is at most $Ct(\alpha+1)^2/X<1/2$.
Enlarging the fixed radius constant and then $C_0$ gives a
contraction. Its fixed point $w$ satisfies
\[
 4tw^3+9-4\pi e^{4w}=-iz,\quad
 |w-q|\leq Ct(\alpha+1)^3/X.
\]
Writing $w=a+ib$ and $Y=-\Ima z\geq0$, real parts give
\[
 4\pi\Rea e^{4w}=4t(a^3-3ab^2)+9+Y\geq cta^3+Y.
\]
We have $|b|<1/2$, and positivity of this real part therefore
places $w$ strictly inside $|b|<\pi/8$.
Also $e^{4a}\asymp X$.

On the disk $|v-w|\leq1/X$, the real theta margin changes
by only $O(1)$, so remains at least $cta^3+ cY$.
The error and its first derivative in~\eqref{eq:slope} are
$O(e^{-4a})$ there. The model derivative equals
\[
 12tw^2-16\pi e^{4w}
        =-4iz+t(12w^2-16w^3)-36;
\]
its norm is between $cX$ and $CX$, because
$ta^3\leq CX/C_0$. Its variation across this disk is $O(1)$.
The Newton map using this fixed model derivative consequently
has derivative $O(1/X)$ on the disk and center displacement
$O(e^{-4a}/X)$. For large $\alpha$ it maps the disk into itself
and is contractive. Its fixed point solves full stationarity.
The retained positive real theta margin and $|\Ima v|<1/2$
keep it in the strip, and its margin exceeds
$c_0t(\Rea v)^3$. It is the asserted admissible actual saddle.
\end{proof}

\begin{proposition}\label{prop:farzeros}
For a sufficiently large fixed constant $C_1$, if $t\geq1$,
$\alpha=\tfrac14\log(|z|/(4\pi))$ is sufficiently large and
\begin{equation}\label{eq:farcondition}
 |z|\geq C_1t(\alpha+1)^3,
\end{equation}
then every zero $H_t(z)=0$ is real and simple.
For each finite $T$, one radius works for every $1\leq t\leq T$.
\end{proposition}
\begin{proof}
Choose $C_1$ at least twice the $C_0$ in the preceding lemma.
If $\Ima z\leq-1$, its actual saddle and
Lemma~\ref{lem:reflection} exclude a zero.
Conjugation excludes $\Ima z\geq1$. In the remaining strip,
put $x=\Rea z$. For large $|z|$ we have
$|x|\geq|z|/2$, and
$\alpha_x=\tfrac14\log(|x|/(4\pi))\leq\alpha$.
Thus $|x|\geq C_0t(\alpha_x+1)^3$, with $\alpha_x$ large.
At this real frequency the far-saddle construction applies.
Proposition~\ref{prop:localactual} proves that the proposed
zero in $|\Ima z|<1$ is real and simple.

For $1\leq t\leq T$, condition~\eqref{eq:farcondition}
and the lower bound on $\alpha$ hold for every sufficiently
large $|z|$, uniformly in $t$, since
$X/(1+\log X)^3\to\infty$. If $T<1$ the assertion is vacuous.
\end{proof}

\section{Whole-plane assembly with one moving center}
The bounded-ratio and far estimates cover all frequencies
only after their scales are compared at the same real center.
\begin{proposition}\label{prop:eventual}
For every sufficiently large $t$, every complex zero of $H_t$
is real and simple.
\end{proposition}
\begin{proof}
Use the actual maximizing center $u$ and
$B=4\pi e^{4u}$ from Lemma~\ref{lem:realcenter}.
For $X=|z|\geq MB$, put
$\alpha=\tfrac14\log(X/(4\pi))=u+d$.
Then $d\geq\tfrac14\log M$ and
\[
 \frac{t(\alpha+1)^3}{X}
 \leq C e^{-4d}\left(1+\frac{d+1}{u}\right)^3
 \leq C(1+d)^3e^{-4d},\qquad u\geq1.
\]
Choose one fixed $M>1$ so large that the last quantity is
at most $1/C_1$ for every $d\geq\tfrac14\log M$.
This is possible by exponential domination.
For late $t$, the far condition and lower bound on $\alpha$
then hold throughout $X\geq MB$.
Proposition~\ref{prop:farzeros} treats that region.
For this same fixed $M$ and the same center $u$,
Proposition~\ref{prop:transition} treats $X\leq MB$.
The two regions cover $\C$. This proves the assertion without
exchanging a fixed-time far limit with a late-time limit.
\end{proof}

\begin{corollary}\label{cor:confine}
There is $R<\infty$ such that every nonreal zero at every
time $t\geq1$ lies in $|z|<R$.
\end{corollary}
\begin{proof}
Choose $T\geq1$ after which Proposition~\ref{prop:eventual}
holds. On $[1,T]$ use the uniform radius from
Proposition~\ref{prop:farzeros}. Beyond $T$ there are no
nonreal zeros. Enlarge the radius to make it positive.
\end{proof}

\section{Attainment and the multiple real zero}
It remains to pass from eventual analytic control to the
literal threshold in~\eqref{eq:starts}. No propagation
property of the quartic flow is used in this passage.

\begin{lemma}\label{lem:closedLP}
If $s_n\to t$ and $\mathsf{LP}(s_n)$ for every $n$, then
$\mathsf{LP}(t)$.
\end{lemma}
\begin{proof}
For each $n$ select one split-real polynomial $Q_n$ within
$1/(n+1)$ of $H_{s_n}$ on $|z|\leq n$. On a fixed compact
set, the error to $H_t$ is eventually bounded by this
quantity plus $\sup|H_{s_n}-H_t|$, which tends to zero by
Proposition~\ref{prop:entire}. Thus $Q_n\to H_t$ locally
uniformly on the whole plane.
\end{proof}

\begin{proof}[Proof of Theorem~\ref{thm:main}]
Proposition~\ref{prop:eventual} and Corollary~\ref{cor:LP}
make $\mathcal S$ nonempty. Proposition~\ref{prop:obstruction}
bounds every start below by a fixed number greater than eight.
Thus $A$ is finite and $A>8$.
For each $t>A$, the definition of an infimum gives some
$T\in\mathcal S$ with $T<t$, whence $\mathsf{LP}(t)$.
Apply Lemma~\ref{lem:closedLP} to $t_n=A+1/(n+1)$ to obtain
$\mathsf{LP}(A)$ as well. Thus $A\in\mathcal S$ and
$\mathcal S=[A,\infty)$. Corollary~\ref{cor:LP} gives
reality of every zero at and after $A$.

Given $\varepsilon>0$, put $T=A-\varepsilon/2$.
Since $T\notin\mathcal S$, some $s\geq T$ fails
$\mathsf{LP}(s)$. Necessarily $s<A$, since all times at or
after $A$ satisfy it. Corollary~\ref{cor:LP} supplies a
nonreal zero at this $s\in(A-\varepsilon,A)$.

Choose $s_n\in(A-1/(n+1),A)$ with nonreal zero $z_n$.
All these times exceed one because $A>8$.
Corollary~\ref{cor:confine} puts all $z_n$ in one compact
disk. After passing to a subsequence, $z_n\to w$ and
$s_n\to A$. Joint continuity gives $H_A(w)=0$, and
threshold reality gives $w\in\R$.

Suppose $c=H_A'(w)\ne0$. Joint derivative continuity provides
a time neighborhood of $A$ and a disk about the real $w$ in
which $|H_s'(q)-c|\leq|c|/2$.
Integration along any segment in that convex disk yields
\[
 |H_s(q)-H_s(p)|\geq\tfrac12|c|\,|q-p|.
\]
For large $n$, both $z_n$ and $\overline z_n$ lie there and
are distinct zeros of the same $H_{s_n}$. This contradicts
the lower bound. Hence $H_A'(w)=0$.
Since $H_A(0)>0$, $w\ne0$. If $w<0$, replace every $z_n$
by $-z_n$ and $w$ by $-w$; evenness preserves all the zero
identities, and its derivative is odd. Thus the limit is a
positive real multiple zero. Its order is at least two and finite, because
$H_A$ is a nonzero entire function. Global eventual
simplicity was already proved in
Proposition~\ref{prop:eventual}.
\end{proof}


\section{Complete divisor polynomials and Gaussian integration}
The polynomial construction used above has a broader classical domain.
Keeping its exact zeros also clarifies what approximation alone can and
cannot justify. Throughout this section let $f$ be an even entire function,
real on the real axis, with $f(0)\ne0$ and
\begin{equation}\label{eq:generalbudgetgrowth}
 |f(z)|\leq C e^{|z|^{3/2}}\qquad(z\in\C).
\end{equation}
Real on the real axis implies $f(\overline z)=\overline{f(z)}$ by the
identity theorem. No reality or simplicity of its zeros is assumed.
For each $R>0$, choose one representative $a$ of every pair $\{a,-a\}$
of zeros in $|a|<R$, and retain its multiplicity $m_a$. Define
\[
 N_R(z)=\prod_{\{a,-a\},\ |a|<R}(1-z^2/a^2)^{m_a}.
\]
This is the complete disk divisor product, not a product that discards
imaginary-axis or nonreal zeros. The choice within each pair has no effect.

\begin{proposition}\label{prop:generaldivisor}
There are zero-free boundary radii $R_n\in(2n+2,2n+3)$ such that the
real even polynomials $P_n=f(0)N_{R_n}$ converge locally uniformly to $f$.
They satisfy $P_n(0)=f(0)$, and their zeros, with multiplicity, are exactly
the zeros of $f$ in $|z|<R_n$. In particular, real coefficients do not
mean that the polynomials split over $\R$.
\end{proposition}
\begin{proof}
There are only countably many excluded zero moduli. In the proof of
Lemma~\ref{lem:polynomial}, the finite disk factorization, even normalized
logarithm and double-zero Schwarz estimate use evenness and
\eqref{eq:generalbudgetgrowth}, with no real-zero assumption. The same
Jensen argument gives $n_f(R)=O(1+R^{3/2})$ for the total number of zeros,
including all complex zeros and multiplicities. The paired denominator
error on $|z|\leq r$ is therefore at most
\[
 \exp\!\left(\frac{n_f(R)r^2}{2R^2}\right)-1\longrightarrow0.
\]
Together with $h_R/h_R(0)\to1$, the exact reconstruction in that proof gives
$f(0)N_R\to f$ along the chosen radii. Every factor has roots exactly
$a,-a$, with the stated multiplicity, so there are no omitted or added
roots. Conjugation permutes the complete zero multiset, hence fixes every
coefficient of $N_R$; those coefficients are real. Its value at zero is
one. These observations prove every assertion.
\end{proof}
The result is a concrete classical factorization consequence, consistent
with Hadamard theory; McMullen's primary course notes state that theory
as Theorem 3.15 and illustrate cancellation of opposite exponential
factors in Theorem 3.18~\cite{mcmullen}. Our preceding disk argument
provides its needed proof here. The corresponding frozen modules are
\code{EvenEntireCanonicalApproximation} and
\code{EvenEntirePolynomialApproximation}.

\subsection*{A common quadratic budget, with an explicit zero strip}
Suppose additionally that every zero satisfies $|\Ima a|\leq b$ for a
fixed $b\geq0$. This hypothesis is retained in the named formal budget
theorem. Write
\[
 S_n=\sum_{\{a,-a\},\ |a|<R_n}\frac{m_a}{a^2}.
\]
Differentiating the finite product at zero gives
$P_n''(0)=-2f(0)S_n$. Cauchy's integral formula on any fixed circle
passes the compact-uniform limit through two derivatives. Thus
\begin{equation}\label{eq:inversezerosum}
 S_n\longrightarrow -\frac{f''(0)}{2f(0)}.
\end{equation}
If $|a|\geq2b$, then $2|\Ima a|\leq|a|$, and direct computation yields
\[
 |a^{-2}|\leq2\Rea(a^{-2}),\qquad
 \Rea(a^{-2})=\frac{(\Rea a)^2-(\Ima a)^2}{|a|^4}.
\]
Indeed the numerator is $|a|^2-2(\Ima a)^2\geq |a|^2/2$.
For $j\geq i$ and $R_i\geq2b$, this gives the absolute annular estimate
\begin{equation}\label{eq:absoluteannulus}
 \sum_{R_i\leq|a|<R_j}\frac{m_a}{|a|^2}
 \leq2\Rea(S_j-S_i)\longrightarrow0
 \quad\hbox{uniformly as }i,j\to\infty.
\end{equation}
Here and below sums use one member of each opposite pair. Our chosen
radii are strictly increasing. Formula~\eqref{eq:inversezerosum} makes
the right side small; the cone inequality also proves its nonnegativity.

\begin{theorem}\label{thm:uniformbudget}
For every $\epsilon>0$ there is $C_\epsilon<\infty$ such that
\[
 |P_n(z)|\leq C_\epsilon e^{\epsilon|z|^2}
 \qquad\hbox{for every }n\geq0\hbox{ and }z\in\C.
\]
The constant is common to all complete divisor approximants, although
it may depend on $f$ and $\epsilon$.
\end{theorem}
\begin{proof}
Choose $i$ so large that every later annulus in
\eqref{eq:absoluteannulus} has sum at most $\epsilon/2$.
Since $|1-w|\leq e^{|w|}$, its product has modulus at most
$\exp(\epsilon|z|^2/2)$. The fixed inner polynomial has a bound
$C e^{\epsilon|z|^2/2}$: each of its finitely many factors satisfies
\[
 |1-cz^2|\leq (1+|c|/\delta)e^{\delta|z|^2},\qquad\delta>0.
\]
To see this, use $1+|c||z|^2\leq
(1+|c|/\delta)(1+\delta|z|^2)$ and $1+x\leq e^x$.
Allocate $\delta=(\epsilon/2)/(1+M)$, where $M$ is the total inner
pair multiplicity. Multiplying its factors proves the inner bound.
Factor $P_n$ into inner and annular products for $n\geq i$, including
its fixed factor $f(0)$. The two budgets add to $\epsilon$.
Each of the finitely many earlier polynomials has its own such bound;
taking their maximum completes the proof for all $n$.
\end{proof}
These are the actual mechanisms in the two frozen modules:
\par\noindent\code{EvenEntireInverseSquareTails};
\par\noindent\code{EvenEntireGaussianBudget}.
\par\noindent The strip is a hypothesis about this
entire function's zeros; it is not an assumption that they are real.

\subsection*{What now passes through an actual Gaussian integral}
Let $Y$ have any fixed real Gaussian probability law, possibly degenerate,
and fix $z,c\in\C$. Then
\begin{equation}\label{eq:gaussdivisorlimit}
 \mathbb E|P_n(z+cY)-f(z+cY)|\longrightarrow0,
 \qquad
 \mathbb E P_n(z+cY)\longrightarrow\mathbb E f(z+cY).
\end{equation}
All these random variables are genuinely integrable. For a direct proof,
choose $a>0$ with $\mathbb E e^{aY^2}<\infty$, available for every real
Gaussian law (and immediate for a constant law). Choose
$\epsilon=a/(4(1+|c|^2))$. The inequality
$|z+cy|^2\leq2|z|^2+2|c|^2y^2$ gives the common majorant
\[
 |P_n(z+cy)|\leq C_\epsilon e^{2\epsilon|z|^2}e^{ay^2}.
\]
The same bound holds for $f$ by pointwise convergence. Twice this
integrable majorant controls the difference, so dominated convergence
proves both limits. This is the fixed-affine Gaussian $L^1$ conclusion
actually mapped in \code{gaussian_L1_limit}; it does not invoke a
formal interchange of an unbounded series of differential operators.

For clarity, a further classical \emph{human deduction} is local
uniformity in the affine parameters: if $|z|\leq A$, $|c|\leq B$, choose
$\epsilon=a/(4(1+B^2))$ once. For each fixed $y$, local uniformity of
$P_n\to f$ gives uniform convergence over that compact parameter set,
because its affine images lie in $|w|\leq A+B|y|$. The supremum of the
continuous error over the compact parameter set is measurable in $y$
(for example use a countable dense subset). The same common majorant
and dominated convergence show that the supremum of the expected
absolute errors tends to zero. This strengthens only the parameter
quantifiers of \eqref{eq:gaussdivisorlimit}, not the class of clocks.

Compact-uniform convergence alone would be insufficient. A precise
human counterexample for standard $Y$ is
\[
 Q_n(z)=1+\frac{z^{2n}}{(2n-1)!!},\qquad n\geq1.
\]
Since $(2n-1)!!\geq n!$, on $|z|\leq A$ its added term is at most
$A^{2n}/n!\to0$ (the ratio tends to zero). Hence $Q_n\to1$ locally
uniformly, with $Q_n(0)=1$, evenness and real coefficients. Nevertheless
$\mathbb E Q_n(Y)=2$ for every $n$, whereas $\mathbb E1=1$; the expected
absolute error is exactly one. These are not asserted to be complete
divisor products of a common function or to have a fixed zero strip.
Their role is to disprove an interchange based solely on compact convergence.

Finally, a Gaussian affine average is a quadratic smoothing operation.
Even when $c$ is complex and the polynomial average can be written as a
finite differential sum, the proof above is an integral domination proof.
It neither identifies a growing quartic deformation with Gaussian
smoothing nor propagates real zeros under that quartic deformation.
The attained threshold and every limitation of the preceding sections
therefore retain exactly their original scope.


\section{A four-point obstruction for the different quadratic flow}
The quadratic comparison can be made more precise without identifying
its threshold with the quartic one. Define the different actual transform
\[
 J_t(z)=\int_0^\infty e^{tu^2}K(u)\cos(uz)\,du,
 \qquad t\in\R,\ z\in\C.
\]
The decay proof of Section 2 gives all absolute integral domains,
entireness and locally uniform time continuity here as well. The same
Young bound $|z|u\leq c|z|^{3/2}+c'u^3$, with the superexponential
theta decay absorbing $tu^2+c'u^3$, gives
$|J_t(z)|\leq C_t e^{|z|^{3/2}}$ after changing constants and splitting
at a fixed large $u$; alternatively allocate the coefficient of the
$|z|^{3/2}$ term to one by choosing its conjugate coefficient suitably.
Also $J_t$ is even, conjugation compatible and $J_t(0)>0$.
The growth class therefore permits the same complete divisor argument.
Only in this section do we use the quadratic time parameter $t$.

\begin{proposition}\label{prop:quadraticnegative}
For some $t<0$, the actual $J_t$ has a nonreal zero.
This conclusion does not specify the negative time or the zero.
\end{proposition}
We give the full four-point proof. It is a classical obstruction
mechanism, not a new quadratic-threshold theorem. Newman already proves
the stronger finite-threshold conclusion in his primary
Theorem 3~\cite{newman}, using his measure classification. The argument
below reconstructs the frozen library's finite-log-factor route instead
of assuming that classification.

\subsection*{The elementary four-point inequality}
For $c\geq0$ let $\ell_c(v)=\log(1+cv)$ on $v\geq0$, and define
\[
 D_c=2\log(1+c/2)-\log(1+c).
\]
For every $y\geq1$,
\begin{equation}\label{eq:fourpointfactor}
 \ell_c(y)\geq \ell_c(1)+(y-1)(\ell_c(2)-\ell_c(1))
                         -2(y-1)^2D_c.
\end{equation}
To verify it, put $u=c/(1+c)$ and $v=y-1$. Both are nonnegative.
The scalar inequality $\log(1+w)\geq w-w^2/2$ for $w\geq0$ follows
by differentiating the difference, whose derivative is $w^2/(1+w)$.
Furthermore
\[
 D_c=\log\frac{(1+c/2)^2}{1+c}
 \geq1-\frac{1+c}{(1+c/2)^2}
 =\left(\frac c{c+2}\right)^2,
 \qquad u\leq\frac{2c}{c+2}.
\]
Here $\log r\geq1-r^{-1}$ for $r>0$ follows from
$\log x\leq x-1$ applied to $x=r^{-1}$. Consequently $u^2\leq4D_c$.
Now $\ell_c(y)-\ell_c(1)=\log(1+uv)$ and
$\ell_c(2)-\ell_c(1)=\log(1+u)\leq u$. Therefore
\[
 \log(1+uv)\geq uv-\tfrac12u^2v^2
                 \geq v\log(1+u)-2v^2D_c,
\]
which proves \eqref{eq:fourpointfactor}.

It follows by summation that every profile
$q(v)=A-Bv+\sum_j\alpha_j\log(1+c_jv)$, with $c_j,\alpha_j\geq0$,
satisfies
\begin{equation}\label{eq:fourpointprofile}
 q(y)\geq q(1)+(y-1)(q(2)-q(1))
       -2(y-1)^2\bigl(2q(1/2)-q(0)-q(1)\bigr),\qquad y\geq1.
\end{equation}
Constant and linear functions satisfy this with equality, so adding
$A-Bv$ changes no curvature term. For a nonzero polynomial $P$ with
all roots real and $P(0)\ne0$, its complete real-root factorization gives
\[
 \log|P(iC\sqrt v)|
 =A_P+\frac12\sum_{r\in\operatorname{roots}(P)}
                    \log\bigl(1+(C/r)^2v\bigr).
\]
All roots are counted with multiplicity; each is nonzero. Hence
$A-Bv+\log|P(iC\sqrt v)|$ satisfies \eqref{eq:fourpointprofile}.
If an entire $f$ in our growth class has only real zeros and $f(0)\ne0$,
Section 11 gives such real-rooted approximants with the exact disk
zeros and preserved origin. At every point on the imaginary axis
$f$ is nonzero: a purely imaginary real zero must be zero itself.
Thus continuity of $\log|\cdot|$ passes their inequalities to
$A-Bv+\log|f(iC\sqrt v)|$ at the five fixed values
$0,1/2,1,2,y$. This step needs no differentiation or infinite logarithm sum.

\subsection*{An actual density smoothing identity}
Let $F(x)=K(|x|)$ and $Y\sim N(0,2)$. This is a bounded, strictly
positive continuous even function. For $\epsilon>0$ put
$S_\epsilon(x)=\mathbb E F(x+\epsilon Y)$. These are genuine positive
integrals. Bounded dominated convergence gives
$S_\epsilon(x)\to F(x)$ as $\epsilon\downarrow0$.
The Gaussian variance is exactly $2\epsilon^2$. Pairing the two
half-lines and completing the square gives
\begin{equation}\label{eq:quadraticdensityheat}
 S_\epsilon(x)=\frac{1}{\sqrt\pi\epsilon}
 e^{-x^2/(4\epsilon^2)}
 J_{-1/(4\epsilon^2)}\left(\frac{ix}{2\epsilon^2}\right).
\end{equation}
Indeed its full-line density is
$(2\sqrt\pi\epsilon)^{-1}e^{-(u-x)^2/(4\epsilon^2)}$.
The sum at $u$ and $-u$, on $u\geq0$, equals
\[
 (\sqrt\pi\epsilon)^{-1}e^{-x^2/(4\epsilon^2)}
 e^{-u^2/(4\epsilon^2)}\cosh(ux/(2\epsilon^2)).
\]
Multiplying by $K(u)$ and using $\cos(iw)=\cosh w$ proves the formula.
The bounded-density Gaussian integral is absolutely integrable;
the half-line cosh integral is absolutely integrable by theta decay.
In particular the rightmost transform value is strictly positive real,
so every logarithm used below is legitimate.

\begin{proof}[Proof of Proposition~\ref{prop:quadraticnegative}]
Suppose instead that $J_t$ has only real zeros for every $t<0$.
For each $\epsilon>0$ take $f=J_{-1/(4\epsilon^2)}$ in the preceding
polynomial-limit inequality. Taking logarithms in
\eqref{eq:quadraticdensityheat} shows that
$q_\epsilon(v)=\log S_\epsilon(\sqrt v)$ is exactly one such profile
with
\[
 A=\log(1/(\sqrt\pi\epsilon)),\qquad
 B=1/(4\epsilon^2),\qquad C=1/(2\epsilon^2).
\]
Consequently \eqref{eq:fourpointprofile} holds for $q_\epsilon$.
Fix $y\geq1$ and send $\epsilon\downarrow0$. Positivity of $F$ and
continuity of the real logarithm imply convergence at all five required
points. Thus the same inequality holds for
$q(v)=\log K(\sqrt v)$.

Its right side is one fixed quadratic polynomial $p(y)=ay^2+by+c$.
For $y\geq1$,
$p(y)\geq-(|a|+|b|+|c|)y^2$. Put $L=|a|+|b|+|c|$.
Then $\log K(x)\geq-Lx^4$ for every $x\geq1$, equivalently
\[
 e^{Lx^4}K(x)\geq1\qquad(x\geq1).
\]
But the actual theta estimate $K(x)\leq C e^{9x-\pi e^{4x}}$
gives $e^{Lx^4}K(x)\to0$: the exponential $e^{4x}$ dominates the
fixed polynomial $Lx^4+9x$. This contradiction proves that at least
one of the strictly negative quadratic times has a nonreal zero.
\end{proof}
The elementary scalar factors, complete polynomial passage and actual
smoothing identity are formalized in
\par\noindent\code{NewmanFourPointBound};
\par\noindent\code{RiemannQuadraticDensitySmoothing};
\par\noindent\code{RiemannQuadraticDensityHeat}.
\par\noindent The frozen endpoint
\code{RiemannQuadraticFiniteThreshold.negative_time_nonreal_exists}
uses the library's finite-threshold route. The direct restriction above
to negative smoothing times is an explicitly labelled human proof of
that same conclusion; it invokes no positive-time strip theorem.

This argument credits established quadratic theory. It gives no
numerical lower bound for a threshold and no statement about $J_0$.
Neither $e^{tu^2}$ in \eqref{eq:quadraticdensityheat} nor the Gaussian
variance can be changed to a quartic weight. It therefore supplies a
precise comparison with the quartic paper, rather than an identification
of the two flows or their eventual-start sets.

\section{Formal correspondence and limits of the conclusion}
The mathematical baseline is the frozen verified commit
\begin{center}\small\ttfamily
bab688bb6140cd6316c80680095a914ea0d346db.
\end{center}
The accompanying source dossier records exact declaration
types, lines, hashes and links, the finite certificate table,
and its independent rational aggregation check. The main
formal namespaces, with prefix
\code{OperatorDerivativeRewrite.Foundation}, are:
\begin{center}
\begin{tabular}{@{}l@{}}
\code{RiemannQuarticUniformFarScale}\\
\code{RiemannQuarticGlobalSimpleZeros}\\
\code{RiemannQuarticActualPolynomialLimit}\\
\code{RiemannQuarticThresholdEndpoint}\\
\code{RiemannQuarticThresholdZeros}\\
\code{RiemannQuarticThresholdCompactness}\\
\code{RiemannQuarticThresholdCollision}.
\end{tabular}
\end{center}
The actual certificate is \code{actual_real_derivative_negative}
in \code{RiemannQuarticLaguerreWitness}.

The human proof groups the actual logarithmic-derivative
estimates into one admissible-saddle calculation and uses
unspecified uniform constants. The formal source retains
explicit constants, model corrections and separate transition
and far budgets. The regrouping is an exposition of the same
final result, not an assertion that every intermediate lemma
above has a new named Lean declaration. The exact final
quantifiers are those of the mapped declarations. The
arithmetic finite inclusions retain their explicit
computer-assisted proofs, while the analytic regrouping and
polynomial reconstruction are supplied as human proofs here.

The threshold is an attained eventual-start threshold.
Neither monotone propagation from each individual reality
time nor failure at every earlier time has been proved here.
The collision argument proves finite multiplicity at least
two, not exactly two or a specified zero trajectory.
The upper threshold is existential, without a numerical
upper value or an optimal quantitative estimate.
No comparison with the classical quadratic threshold or
implication for the Riemann hypothesis is supplied.

De Bruijn~\cite{debruijn} and Newman~\cite{newman} are primary
antecedents for real-zero Fourier transforms and quadratic
propagation. Their exact deformation, sign and quantifiers
must be retained in any priority comparison. Preliminary
comparison does not establish novelty for the combined
quartic package; novelty remains unresolved.

\begin{thebibliography}{9}
\bibitem{debruijn} N. G. de Bruijn,
The roots of trigonometric integrals,
\emph{Duke Math. J.} \textbf{17} (1950), 197--226.
\href{https://doi.org/10.1215/S0012-7094-50-01720-0}
{doi:10.1215/S0012-7094-50-01720-0}.
\bibitem{newman} C. M. Newman,
Fourier transforms with only real zeros,
\emph{Proc. Amer. Math. Soc.} \textbf{61} (1976), 245--251.
\url{https://sites.math.northwestern.edu/~auffing/papers/Newman.pdf}.
\bibitem{dlmf} NIST Digital Library of Mathematical Functions,
\S20.7(viii), equation 20.7.32, modular theta identity.
\url{https://dlmf.nist.gov/20.7.E32}.
\bibitem{mcmullen} C. McMullen,
\emph{Advanced Complex Analysis}, Harvard Math 213a course notes,
November 30, 2023, Theorems 3.15 and 3.18, printed pages 80--82.
\url{https://people.math.harvard.edu/~ctm/home/text/class/harvard/213a/23/html/home/course/course.pdf}.
\end{thebibliography}
\end{document}
